% ---------------------------------------------------------------------------
% Capítulo 2 --- Marco teórico
% ---------------------------------------------------------------------------
% Redactado en tercera persona (norma 2.4). Los acrónimos ya introducidos en
% la Introducción no se repiten en extenso salvo cuando aclara la exposición.
% Decimales con coma y unidades del Sistema Internacional. Mínimo dos páginas
% por capítulo (norma 2.6).

\chapter{Marco teórico}

Este capítulo reúne los fundamentos sobre los que se apoya el desarrollo. Avanza desde el
problema que motiva el proyecto ---la fatiga de falsos positivos en el triaje de alertas--- hacia las
plataformas que hoy lo abordan, los modelos de lenguaje que este trabajo incorpora, la técnica de
recuperación que corrige su principal debilidad, los riesgos de aplicarlos a este dominio, el
razonamiento que sostiene que la decisión final no debe tomarla el modelo, y por último los
criterios con que se mide el resultado.

% ===========================================================================
\section{El problema del triaje de alertas de seguridad}
\label{sec:problema-triaje-referencia}
% ===========================================================================

Un sistema de detección basado en reglas y firmas dispara una alerta cuando una secuencia de
eventos coincide con un patrón previamente definido: una autenticación fallida, un puerto que
empieza a recibir tráfico, un proceso que no debería ejecutarse. Esa coincidencia es determinista y
auditable ---se puede señalar exactamente qué regla se disparó---, pero no incorpora ningún juicio
sobre si, en el contexto del activo concreto que la generó, esa coincidencia representa una
amenaza real. La misma regla se dispara igual sobre un servidor expuesto a internet que sobre uno
que no lo está, y el resultado es un volumen de falsos positivos que crece con el tamaño de la red
sin que crezca en la misma proporción el riesgo real.

El \emph{triaje} es la tarea de decidir, para cada alerta, si merece la atención de un analista y
con qué prioridad. Cuando esa tarea se realiza manualmente sobre un volumen de alertas que supera
la capacidad del equipo, dos fallos son posibles y ambos son costosos: tratar como amenaza algo que
no lo es, que consume tiempo del analista sin reducir el riesgo, y descartar como ruido algo que sí
lo es, que es exactamente el fallo que un sistema de seguridad no puede permitirse. Las evaluaciones
independientes de las plataformas comerciales de detección más recientes miden de forma explícita
el volumen de alertas que un analista debe procesar como un criterio de comparación entre
proveedores \cite{mitreattackevals}, lo que confirma que este problema no es una limitación de las
herramientas más antiguas, sino una característica persistente del propio enfoque basado en reglas,
incluso en sus implementaciones más avanzadas.

% ===========================================================================
\section{Detección y respuesta: EDR, XDR y MDR}
% ===========================================================================

Tres siglas dominan el mercado de plataformas de detección y conviene distinguirlas, porque operan
en planos distintos.

\textbf{EDR} (\emph{Endpoint Detection and Response}, detección y respuesta en el \emph{endpoint})
es la tecnología que recoge telemetría de un equipo concreto ---procesos, ficheros, registro,
conexiones de red locales--- y responde sobre ese mismo equipo: aislarlo, terminar un proceso,
ponerlo en cuarentena.

\textbf{XDR} (\emph{Extended Detection and Response}, detección y respuesta extendida) correlaciona
telemetría de varias capas ---\emph{endpoint}, red, identidad, correo, nube--- en una sola plataforma
analítica. El EDR es, en este esquema, una de las fuentes que alimentan al XDR; lo que este añade
es la correlación entre fuentes que, por separado, solo verían una parte del incidente.

\textbf{MDR} (\emph{Managed Detection and Response}, detección y respuesta gestionada) no es una
tecnología sino un servicio: un equipo humano que opera un EDR o un XDR por cuenta del cliente,
haciendo el triaje, la investigación y la respuesta que el cliente no tiene capacidad de operar por
sí mismo. Un MDR que opera sobre un XDR se denomina en ocasiones MXDR (\emph{Managed XDR}).

Las plataformas comerciales líderes despliegan típicamente un agente propietario a nivel de
sistema operativo, más rico que las fuentes de telemetría de código abierto ---como Sysmon, osquery o
Zeek--- y con capacidad de respuesta activa integrada. Esas fuentes abiertas, sin embargo, no quedan
excluidas del ecosistema: numerosas plataformas las ingieren como fuente complementaria en su capa
de agregación, y los servicios MDR que no exigen un agente comercial concreto ---el modelo que en
este trabajo se denomina telemetría-agnóstico--- operan directamente sobre la telemetría que el
cliente ya tenga. Este proyecto se inserta precisamente en ese segundo modelo: no sustituye al
agente ni a la fuente de telemetría, sino que añade la capa de clasificación, priorización y
justificación explicable sobre las alertas que esa fuente ya produce.

% ===========================================================================
\section{MITRE ATT\&CK como taxonomía de referencia}
% ===========================================================================

MITRE ATT\&CK (\emph{Adversarial Tactics, Techniques, and Common Knowledge}) es una matriz pública
y de acceso libre que cataloga las tácticas y técnicas que un atacante emplea a lo largo de un
incidente, desde el acceso inicial hasta la exfiltración, con un identificador único por técnica
---por ejemplo, T1110 para el acceso por fuerza bruta, con la subtécnica T1110.001 para el intento
por diccionario de contraseñas \cite{mitreattack}. Su valor para este trabajo es doble: sirve como
vocabulario común para expresar la hipótesis de ataque dentro de la justificación de cada decisión,
y sirve como criterio de comparación en las evaluaciones independientes de plataformas comerciales
citadas en la sección anterior. MITRE mantiene además ATLAS, una matriz equivalente centrada en las
amenazas contra los propios sistemas de inteligencia artificial \cite{mitreatlas}, pertinente en la
medida en que este trabajo introduce uno de esos sistemas en el propio camino de decisión.

% ===========================================================================
\section{Modelos de lenguaje de gran escala aplicados a la seguridad}
% ===========================================================================

Un modelo de lenguaje de gran escala (LLM) es una red neuronal entrenada para predecir la
continuación más probable de una secuencia de texto, sobre un volumen de parámetros que va de
varios cientos de millones a varios cientos de miles de millones. Aplicado a la seguridad, conviene
distinguir tres formas de uso, que no compiten entre sí sino que resuelven partes distintas del
problema de triaje:

\begin{itemize}
  \item \textbf{Codificador clasificador} (\emph{encoder}, con arquitectura de tipo BERT y del
        orden de cien millones de parámetros): asigna una clase y una confianza numérica a una
        entrada, tras un ajuste fino sobre datos etiquetados. Su coste de ejecución es muy bajo
        ---del orden de milisegundos sobre un procesador convencional--- y es determinista, pero no
        redacta ninguna explicación y no generaliza bien fuera de lo que vio durante el
        entrenamiento.
  \item \textbf{Modelo generativo instruido} (de varios miles de millones de parámetros): redacta
        el razonamiento, resume un incidente o traduce una consulta en lenguaje natural. Su coste es
        alto ---segundos y varios gigabytes de memoria--- y no es fiable como clasificador numérico ni
        determinista por defecto.
  \item \textbf{Modelo especializado de dominio}: un codificador o un generativo con
        preentrenamiento continuado sobre corpus de ciberseguridad. \textbf{SecureBERT} demostró que
        esa especialización mejora las tareas de clasificación del dominio frente al modelo genérico
        equivalente \cite{aghaei2022securebert}, y \textbf{Foundation-Sec-8B}, un modelo generativo
        de ocho mil millones de parámetros con preentrenamiento continuado sobre corpus de
        inteligencia de amenazas, declara explícitamente el triaje de un centro de operaciones de
        seguridad como uno de sus casos de uso objetivo \cite{foundationsec8b}. La especialización
        aporta vocabulario ---el modelo reconoce sin explicación previa un identificador de
        vulnerabilidad o una técnica ATT\&CK--- pero no aporta garantías: sigue pudiendo afirmar con
        aparente seguridad algo que la entrada no contiene.
\end{itemize}

La conclusión que este marco teórico traslada al diseño del prototipo es que clasificar y
justificar son tareas de naturaleza distinta y conviene resolverlas con componentes distintos: un
codificador da una clase estable, barata y con una confianza calibrable; un modelo generativo
produce el texto que una persona necesita leer para decidir. Pedirle ambas cosas al mismo
componente tiende a producir puntuaciones de confianza inventadas con apariencia de cálculo.

% ===========================================================================
\section{Generación aumentada por recuperación}
% ===========================================================================

La generación aumentada por recuperación (RAG) combina un modelo de lenguaje con un mecanismo de
recuperación de información sobre un corpus documental externo: ante una consulta, el sistema
recupera primero los fragmentos de texto más relevantes de ese corpus y se los entrega al modelo
como parte de su contexto, antes de que este redacte la respuesta \cite{lewis2020rag}. Dos
consecuencias hacen de esta técnica algo más que un atajo de ingeniería. La primera es que el
conocimiento deja de residir exclusivamente en los parámetros del modelo ---fijados en el momento de
su entrenamiento--- y pasa a un índice que puede actualizarse sin reentrenar nada. La segunda es que
la respuesta queda \emph{anclada}: cada afirmación puede, en principio, rastrearse hasta el
fragmento del corpus que la sustenta, lo que reduce ---sin eliminar--- el riesgo de que el modelo
invente un dato plausible.

La forma más simple de recuperación ---la empleada en este trabajo--- representa cada documento del
corpus y cada consulta como un vector numérico mediante un modelo de \emph{embeddings}, y ordena
los documentos por su similitud coseno con la consulta:

\begin{equation}
  \cos(\mathbf{q}, \mathbf{d}) = \frac{\mathbf{q} \cdot \mathbf{d}}{\lVert \mathbf{q} \rVert \, \lVert \mathbf{d} \rVert}
\end{equation}

donde $\mathbf{q}$ es el vector de la consulta y $\mathbf{d}$ el de un documento del corpus; se
seleccionan los $k$ documentos de mayor similitud, con $k$ típicamente pequeño (en este trabajo,
$k=3$) para no saturar el contexto del modelo con pasajes de relevancia marginal. Existen técnicas
de recuperación más sofisticadas ---dispersas, como BM25, o densas entrenadas específicamente para la
tarea de recuperación, como en los sistemas de tipo \emph{dense passage retrieval}--- así como
arquitecturas de recuperación por grafo; este trabajo emplea la variante más simple posible, por
similitud de \emph{embeddings} genéricos, dado que el objetivo no es optimizar la recuperación en sí
sino verificar si corregir el conocimiento del modelo mejora la corrección de su justificación.

Un riesgo documentado de esta técnica, especialmente relevante cuando la consulta y el corpus se
formulan en idiomas o registros distintos, es que la utilidad de un fragmento recuperado no es
uniforme según su posición en el contexto que finalmente recibe el modelo: los fragmentos situados
en el centro de un contexto largo se aprovechan peor que los situados en sus extremos
\cite{liu2023lost}. Esta observación motivó, en el diseño del justificador de este trabajo,
mantener el número de pasajes recuperados deliberadamente bajo.

% ===========================================================================
\section{Agentes con herramientas y el patrón de razonar y actuar}
\label{sec:agentes}
% ===========================================================================

Un \emph{agente} basado en un modelo de lenguaje es un modelo situado dentro de un bucle: en lugar
de producir una única respuesta, alterna pasos de razonamiento con la invocación de herramientas
externas, observa lo que estas devuelven y decide con esa información si continuar o concluir. La
formulación que popularizó este esquema es el patrón conocido como \emph{ReAct}, contracción de
\emph{reasoning and acting}, que entrelaza explícitamente las trazas de razonamiento con las
acciones sobre el entorno y muestra que ambas se refuerzan: el razonamiento guía qué herramienta
usar, y lo observado corrige el razonamiento siguiente \cite{yao2023react}.

Este trabajo emplea el patrón en dos puntos distintos, y conviene separarlos porque su riesgo no es
el mismo. El primero es la \textbf{recuperación agéntica}: en vez de construir la consulta al corpus
con una plantilla fija sobre los campos de la alerta, es el propio modelo el que formula qué
necesita consultar antes de justificar. El segundo es el \textbf{agente de mitigación}, que ante un
incidente cuya contención en el equipo afectado no basta ---porque ese equipo está fuera de
alcance--- razona la escalada hacia el siguiente dispositivo con capacidad de contención. Este
segundo uso no procede de la literatura sino de una necesidad operativa señalada por el tutor
industrial de la empresa: en una red en producción el equipo atacado puede estar caído,
comprometido o sin canal de administración, y quedarse sin capacidad de contención por ello no es
aceptable.

La tensión con el resto de este marco teórico es evidente y merece enunciarse en lugar de
disimularse: un agente que planifica libremente y encadena herramientas por su cuenta es
precisamente lo que la exigencia de auditabilidad de la sección \ref{sec:politica-perfil} rechaza.
La forma de conservar la utilidad del patrón sin heredar ese problema es acotarlo por construcción,
y en este trabajo se acota de tres maneras simultáneas: el agente \textbf{selecciona} dentro de un
catálogo cerrado de acciones y nunca inventa una nueva; cada paso de la escalada requiere
aprobación humana explícita antes de ejecutarse; y la decisión de contener ---el qué--- sigue
siendo determinista, de modo que el agente elige la estrategia, no si hay que actuar. Un agente así
no es autónomo en el sentido fuerte del término: es un selector razonado sobre un espacio de
acciones cerrado y supervisado.

% ===========================================================================
\section{Riesgos de aplicar un modelo de lenguaje al triaje de seguridad}
% ===========================================================================

Introducir un modelo de lenguaje en el camino de decisión de una alerta de seguridad tiene riesgos
específicos de este dominio, más allá de los riesgos generales de cualquier aplicación basada en
LLM catalogados por el proyecto abierto de seguridad en aplicaciones web \cite{owaspllmtop10} y por
la evaluación de seguridad de modelos de Meta conocida como CyberSecEval \cite{purplellama}.

\textbf{Inyección de instrucciones desde campos controlados por el atacante.} El contenido de un
registro de eventos ---el nombre de usuario probado, la cadena de \emph{user-agent}, el mensaje
completo del evento--- lo escribe, en un intento de intrusión, la propia parte atacante. Si esos
campos se concatenan sin distinguirlos de las instrucciones del sistema, el atacante dispone de un
canal directo hacia el componente que analiza su propio ataque. La respuesta de diseño es tratar
siempre esos campos como datos delimitados, nunca como instrucciones.

\textbf{Alucinación con apariencia de fundamento.} El fallo más peligroso no es el error evidente,
sino la explicación bien redactada que cita un detalle que la entrada no contiene; un analista con
prisa la acepta precisamente porque suena a análisis correcto. El contrapeso es exigir un anclaje
verificable: toda afirmación debe poder rastrearse a un campo concreto de la alerta o del contexto
recuperado, y esa verificación puede automatizarse.

\textbf{Ausencia de determinismo.} La misma entrada puede producir salidas distintas en dos
ejecuciones sucesivas de un modelo generativo. Para la operación es una incomodidad menor; para la
evaluación es descalificante, porque impide que los resultados sean reproducibles. Se controla
fijando una temperatura de muestreo próxima a cero y registrando la versión exacta del modelo y de
las instrucciones empleadas en cada ejecución.

\textbf{Corte de conocimiento.} Un modelo no sabe nada de lo publicado después de su entrenamiento,
así que cualquier afirmación suya sobre el panorama de amenazas actual es sospechosa por
construcción; el conocimiento reciente debe entrar por el contexto que se le suministra, no salir
de sus propios parámetros.

\textbf{Superficie de fallo adicional.} Un componente más que puede fallar, colgarse o rendir por
debajo de lo esperado exige un mecanismo de degradación: si el modelo no responde o su respuesta no
pasa la verificación de anclaje, la decisión debe seguir un camino alternativo ---una plantilla
determinista, en este trabajo--- y nunca descartarse en silencio.

% ===========================================================================
\section{Reglas y firmas frente a análisis asistido por modelos de lenguaje}
% ===========================================================================

La tabla \ref{tab:reglas-vs-ia} resume la comparación entre ambos enfoques sobre los atributos que
importan para el triaje.

\begin{table}[H]
  \centering
  \caption{Reglas y firmas frente a análisis asistido por modelos de lenguaje}
  \label{tab:reglas-vs-ia}
  \begin{tabular}{@{}p{0.24\textwidth}p{0.34\textwidth}p{0.34\textwidth}@{}}
    \toprule
    & \textbf{Reglas y firmas} & \textbf{Análisis asistido por modelo} \\
    \midrule
    Decisión & Determinista y repetible & Probabilístico; repetible solo si se fuerza \\
    Auditabilidad & Total: se lee la regla que se disparó & Indirecta: se lee la justificación, no
      el cálculo \\
    Contexto del activo & Ninguno & Su mejor aportación, si se le suministra \\
    Explicación & El nombre de la regla & Razonamiento en lenguaje natural \\
    Generalización & Nula ante lo no previsto & Parcial: reconoce variantes de lo ya visto \\
    Fallo típico & Falso positivo por falta de contexto & Afirmación segura y equivocada \\
    \bottomrule
  \end{tabular}
\end{table}

Ninguno de los dos enfoques sustituye por completo al otro: las reglas son buenas detectando y
malas decidiendo qué importa, y el modelo es débil detectando por sí solo pero fuerte explicando y
priorizando lo que ya fue detectado. El diseño que este trabajo adopta usa las reglas como fuente
---de ahí que Wazuh se conserve como origen de las alertas--- y el modelo como capa de triaje sobre
ellas, nunca al revés.

% ===========================================================================
\section{Políticas de decisión determinista y separación entre analizar y actuar}
\label{sec:politica-perfil}
% ===========================================================================

Un motor de triaje que, además de clasificar una alerta, decidiera con el mismo componente
generativo qué acción de respuesta ejecutar, heredaría todos los riesgos de la sección anterior en
el punto de mayor consecuencia del sistema: el momento de actuar sobre la infraestructura del
cliente. La alternativa que sostiene el diseño de este trabajo separa esa decisión en dos capas con
motivos de cambio distintos: una \emph{política} determinista, expresada como una tabla indexada
por la clase de la alerta y el contexto del activo, que \emph{propone} la acción de menor impacto
capaz de resolver la amenaza; y un \emph{perfil de cliente} configurable, que \emph{filtra} esa
propuesta según lo que el negocio del cliente concreto permite, degradándola a una de menor impacto
o reteniéndola para validación humana. El modelo de lenguaje interviene únicamente para redactar la
justificación de una decisión ya tomada por la política, nunca para tomarla. Esta separación es, en
sí misma, una aplicación directa del principio de auditabilidad: un error se puede atribuir a la
clasificación, a la política o al perfil, y se sabe cuál de los tres corregir, algo que no sería
posible si toda la decisión residiera dentro de un único componente generativo.

Esta separación es también la que permite incorporar el agente descrito en la sección
\ref{sec:agentes} sin renunciar a la auditabilidad. El agente no sustituye a ninguna de las tres
capas: actúa dentro de la tercera, eligiendo entre las acciones que el catálogo ofrece y que el
perfil ya ha autorizado, cuando la contención en el equipo afectado resulta insuficiente. La
frontera se mantiene, por tanto, en el mismo sitio: la política determina \emph{qué} debe ocurrir y
el modelo, en su papel de agente, propone \emph{cómo} lograrlo dentro de lo permitido, sujeto a
aprobación humana en cada paso.

% ===========================================================================
\section{Métricas de evaluación de un clasificador binario}
% ===========================================================================

Evaluar si el triaje asistido mejora sobre el método de referencia exige medir su desempeño como
clasificador binario ---amenaza o no amenaza--- sobre una matriz de confusión con cuatro categorías:
verdadero positivo ($VP$), cuando el sistema marca como amenaza una alerta que efectivamente lo es;
falso positivo ($FP$), cuando la marca como amenaza sin serlo; verdadero negativo ($VN$), cuando
correctamente no la marca; y falso negativo ($FN$), cuando una amenaza real pasa desapercibida. De
esa matriz se derivan las métricas de la tabla \ref{tab:metricas-clasificacion}.

\begin{table}[H]
  \centering
  \caption{Métricas de clasificación empleadas en la evaluación}
  \label{tab:metricas-clasificacion}
  \begin{tabular}{@{}p{0.22\textwidth}p{0.28\textwidth}p{0.42\textwidth}@{}}
    \toprule
    \textbf{Métrica} & \textbf{Fórmula} & \textbf{Qué mide} \\
    \midrule
    Precisión ($P$) & $\dfrac{VP}{VP+FP}$ & De lo marcado como amenaza, cuánto lo era realmente \\
    Exhaustividad ($R$) & $\dfrac{VP}{VP+FN}$ & De las amenazas reales, cuántas se detectaron \\
    Medida $F_1$ & $\dfrac{2\,P\,R}{P+R}$ & Equilibrio entre precisión y exhaustividad \\
    Tasa de falsos positivos & $\dfrac{FP}{FP+VN}$ & El ruido operativo que sufre el analista \\
    \bottomrule
  \end{tabular}
\end{table}

De estas cuatro, la tasa de falsos positivos es la que responde más directamente al problema
planteado en la sección \ref{sec:problema-triaje-referencia}: mide exactamente la proporción de
ruido que un analista tendría que descartar manualmente. La exhaustividad, por su parte, es la
métrica que no puede sacrificarse bajo ninguna circunstancia: un sistema que reduce el ruido a costa
de dejar pasar amenazas reales no es una mejora, sino un riesgo nuevo.

% ===========================================================================
\section{Entornos de red virtualizados para experimentación en seguridad}
% ===========================================================================

Ejercitar y medir un motor de triaje exige un entorno donde generar actividad legítima y hostil
sobre una red con vulnerabilidades conocidas, sin exponer una infraestructura productiva. Las
plataformas de este tipo se dividen en dos familias con una diferencia determinante para este
trabajo: las que \emph{simulan} un dispositivo ---reproducen su comportamiento observable sin
ejecutar su sistema operativo real--- y las que lo \emph{virtualizan} o lo \emph{emulan}, ejecutando
una imagen real del sistema dentro de un contenedor o una máquina virtual. Solo la segunda familia
expone servicios genuinos que una herramienta de escaneo de vulnerabilidades pueda encontrar, lo
que descarta a los simuladores puramente didácticos como plataforma de laboratorio para este
proyecto.

Containerlab es una de las plataformas de esa segunda familia: define la topología de red como un
fichero de texto versionable, orquesta contenedores y máquinas virtuales ligeras según ese fichero,
y crea automáticamente una red de gestión fuera de banda que conecta todos los nodos entre sí
\cite{containerlabdocs}. Su reproducibilidad ---una topología descrita en un fichero puede destruirse
y recrearse de forma idéntica--- es lo que hace posible que las condiciones de una campaña de
evaluación puedan repetirse de manera fiel, un requisito directo de la reproducibilidad exigida a
la propia evaluación del capítulo \ref{cap:metodologia}.
